\documentclass[10pt,a4paper]{amsart}
\usepackage[margin=19mm]{geometry}
\usepackage{amsmath,amssymb,mathtools}
\usepackage[scr=rsfs,cal=euler]{mathalfa}
\usepackage{xfrac}
\usepackage[unicode,hidelinks,bookmarks=false]{hyperref}
\newcommand{\ie}{i.e.\ }
\newcommand{\cf}{cf.\ }
\newcommand{\resp}{respectively\ }
\newcommand{\Subsection}{\subsection}
\providecommand{\nobreakdash}{\nobreak\hbox{-}}
\setlength{\emergencystretch}{2em}
\multlinegap0pt
\usepackage{amssymb}
\usepackage{enumerate}
\usepackage{xcolor}
\newtheorem{lem}{Lemma}
\def\P{\mathbf{P}}
\title{Galton-Watson processes, simple varieties of trees and Khinchin families}
\author{V\'{\i}ctor J. Maci\'a}
\date{Source: arXiv:2506.18370v7 (CC BY 4.0)}
\begin{document}
\maketitle

\noindent\emph{Source and licence.} Galton-Watson processes, simple varieties of trees and Khinchin families. Version arXiv:2506.18370v7 (CC BY 4.0), distributed by arXiv under the Creative Commons Attribution 4.0 International licence, http://creativecommons.org/licenses/by/4.0/, as recorded in the arXiv licence metadata for this exact version and shown on the abstract page https://arxiv.org/abs/2506.18370v7. Full licence notice: \texttt{licenses/arxiv-2506.18370-CC-BY-4.0.txt}.\par\smallskip

\noindent\textit{Editorial scope.} Counted unit: the article's lem weight\_prob (source0.tex lines 641--647). The article's complete original proof of it is delivered with it (source0.tex lines 649--665, one \texttt{proof} environment of 17 lines). No mathematics is written, repaired or re-derived: the statement and the proof are the article's own lines, in the article's own notation, and no retained line is substituted or re-flowed.  No further source-local context is needed: every cross-reference of every retained line resolves inside the two delivered spans.  Nothing was omitted from any retained span, so the package carries no omission record.  No retained line contains a citation, so the wrapper carries no bibliography and no bibliography entry.  No retained line contains a figure environment, an image-inclusion command or drawing code, so no image is delivered and the package declares no asset dependency.  Editorial formatting only, in addition: an AMS \texttt{amsart} preamble that restates the article's own declarations and macros used by the retained lines, namely the environments \texttt{lem} on separate counters, together with the standard packages those lines need.  The retained environments print in this document as Lemma 1 in order of appearance, whereas the article prints the same items with the numbering of its own full document.  The complete native source of the article as distributed in the arXiv e-print for this exact version is bundled unchanged as \texttt{sources/180326/source0.tex}.
	\begin{lem} \label{lemma: weight_prob}
		Fix \( t \in (0,R_{\psi}) \). Then, for any \( n \geq 1 \) and \( a \in \mathcal{G}_n \), we have
		\begin{align*}
			\P(T_t = a) = \frac{\omega(a)t^{n-1}}{\psi(t)^n}\,,  
		\end{align*}
		with the convention that $0^0 = 1$. Here, \( \omega(a) \) denotes the weight of the rooted plane tree \( a \in \mathcal{G}_n \).
	\end{lem}

	\begin{proof} 
		For any \( n \geq 1 \) and \( a \in \mathcal{G}_n \), we have
		\begin{align*}
			\P(T_t = a) &= \prod_{j = 0}^{n-1} \P(Y_t = j)^{k_j(a)} 
			= \prod_{j = 0}^{n-1} \left( \frac{b_j t^j}{\psi(t)} \right)^{k_j(a)}.
		\end{align*}
		Recalling that for a tree \( a \in \mathcal{G}_n \), the following relations hold:
		\begin{align*}
			k_0(a) + k_1(a) + k_2(a) + \dots + k_{n-1}(a) &= n\,, \\
			k_1(a) + 2k_2(a) + \dots + (n-1)k_{n-1}(a) &= n-1\,,
		\end{align*}
		these equations express that the tree has \( n \) nodes and \( n-1 \) edges (or, equivalently, $n-1$ nodes that descend from a parent node, we don't count the root), we obtain that
		\begin{align*}
			\P(T_t = a) &= \prod_{j = 0}^{n-1} b_j^{k_j(a)} \frac{t^{n-1}}{\psi(t)^n} 
			= \frac{\omega(a)t^{n-1}}{\psi(t)^n}.
		\end{align*}
	\end{proof}

\end{document}
